
Die digitale Transformation ist für Unternehmen zunehmend zu einer strategischen Notwendigkeit geworden. 
Unabhängig von ihrer Größe sehen sich Unternehmen mit der Herausforderung konfrontiert, Geschäftsprozesse 
zu digitalisieren, Informationen schneller verfügbar zu machen und auf veränderte Marktanforderungen zu 
reagieren. Für KMU entsteht hierbei jedoch ein besonderes Spannungsfeld: Einerseits steigt der Druck zur 
Digitalisierung kontinuierlich an, andererseits stehen für entsprechende Transformationsvorhaben häufig 
nur begrenzte Ressourcen zur Verfügung \parencite{Leyh2015,Leyh2019}.

Der steigende Wettbewerbsdruck zwingt Unternehmen dazu, ihre Prozesse kontinuierlich effizienter zu gestalten 
und Kosten zu reduzieren, um langfristig wettbewerbsfähig zu bleiben \parencite{Leyh2015}. Gleichzeitig 
verändern sich die Erwartungen von Kunden zunehmend. Digitale Dienstleistungen, kurze Reaktionszeiten 
sowie transparente Informationsbereitstellung werden immer stärker vorausgesetzt \parencite{Leyh2019}. 
Darüber hinaus führt die Digitalisierung zu höheren Anforderungen an die Vernetzung betrieblicher 
Prozesse und Daten. Informationen müssen abteilungsübergreifend verfügbar sein und unterschiedliche 
Unternehmensbereiche enger miteinander zusammenarbeiten, um effiziente und durchgängige Wertschöpfungsprozesse 
zu ermöglichen \parencite{Leyh2019}.

Dem steigenden Digitalisierungsdruck stehen in KMU häufig begrenzte Ressourcen gegenüber. Insbesondere 
finanzielle Restriktionen erschweren die Umsetzung größerer Digitalisierungs- und IT-Projekte \parencite{Faruque2024}. 
Gleichzeitig verfügen viele KMU nur über geringe interne IT-Kapazitäten und beschäftigen vergleichsweise 
wenige spezialisierte Fachkräfte \parencite{Xia2009, Leyh2015}. Hinzu kommen Wissens- und Erfahrungsdefizite im 
Umgang mit komplexen IT-Systemen und Digitalisierungsprojekten, wodurch strategische Entscheidungen sowie 
die Planung entsprechender Vorhaben zusätzlich erschwert werden \parencite{Leyh2019}.

Aus diesem Spannungsfeld ergibt sich die Frage, wie KMU digitale Transformationsprozesse trotz begrenzter 
Ressourcen erfolgreich gestalten können. ERP-Systeme werden in der Literatur häufig als mögliche 
technologische Grundlage betrachtet, da sie Prozesse, Daten und Informationsflüsse innerhalb eines 
Unternehmens integrieren können. Gleichzeitig stellt sich die Frage, ob und unter welchen Bedingungen 
ERP-Systeme unter den spezifischen Rahmenbedingungen von KMU tatsächlich wirksam werden können.
